\section{Three-way comparison over the energy range}
\label{sec:comp-independent-total-comparison}
The coefficient compared here is
\begin{equation}
 C(s)=C_1(s-1)=\frac{m^2\sigma_1}{\alpha_{\mathrm T}^3},\qquad
 s=\frac{s_{\mathrm{phys}}}{m^2}=1+\frac{2\omega}{m}.
\end{equation}
The photon energy $\omega$ is measured in the initial electron rest
frame. The seven energies span $10^{-4}\leq\omega/m\leq10^5$.
The first value at each energy in
Table~\ref{tab:comp-independent-total-50} is our independent analytic
coefficient, freshly evaluated from the collected expression in
Equation~\eqref{eq:comp-total-energy-collected}. Its 50 displayed digits
also agree with the original assembled sector primitives.
All 50 displayed significant digits agree between two runs at
65 and 85 requested decimal digits, each with 40 guard digits.
This agreement checks numerical stability. The endpoint series used
to evaluate our primitives have truncation-error estimates, but a
rigorous bound on every omitted tail has not been established.
The displayed digits therefore do not constitute a certified
50-digit error bound.

The two subsequent entries are absolute coefficient differences,
\begin{equation}
 \Delta C_{\mathrm{num}}=C_{\mathrm{num}}-C(s),\qquad
 \Delta C_{\mathrm{ref}}=C_{\mathrm{ref}}-C(s).
\end{equation}
The differences use the full saved analytic value before rounding it
to the 50 digits displayed in the table.
For $C_{\mathrm{num}}$, the virtual and soft terms are integrated
numerically over the elastic angle from the pointwise amplitude.
The real term is a separate four-dimensional phase-space integral
of the six-diagram double-Compton matrix element, with its local
soft subtraction. Neither calculation calls the analytic total
or its integrated GPL primitives.

The real quadrature error estimate includes both the integrator's
reported error and the change upon refinement. The virtual estimate
includes angular refinement and an arithmetic accuracy floor.
We add these estimates rather than treating them as independent
statistical errors. Thus the $\pm$ entries are numerical error
estimates, not confidence intervals or experimental uncertainties.
The numerical integration is not required to reproduce 50 digits.

The reference coefficient is Equation (6) of Lee, Schwartz and
Zhang~\cite{qed-processes-lsz}. It is evaluated separately, after
our analytic result is fixed. Its error estimate uses two working
precisions, changes in the special-function integration endpoint,
an analytic bound on the omitted tail, and a rounding scale based
on the sizes of its terms. It does not use the difference from our
answer to estimate its own error. This comparison tests numerical evaluation of the two analytic
expressions. Section~\ref{sec:compton-analytic-equivalence} separately
establishes their exact equality by comparing every coefficient and
fixing their common physical continuation.

\begin{table}[!htbp]
\centering
\comptonprecisiontable
\begin{tabular}{rll}
\toprule
$\omega/m$ & Calculation & Coefficient or difference\\
\midrule
$0.000100$ & Our analytic & $0.00000062423288400014031134169341516503703552376453828153$ \\
 & Numerical difference & $5.30\times10^{-15}\;\pm\;2.0\times10^{-12}$ \\
 & Reference difference & $-4.03\times10^{-80}\;\pm\;5.6\times10^{-58}$ \\
\addlinespace[3pt]
$0.0100$ & Our analytic & $0.0029187564665568768995602481244244431628302067748237$ \\
 & Numerical difference & $1.11\times10^{-10}\;\pm\;2.0\times10^{-8}$ \\
 & Reference difference & $2.37\times10^{-86}\;\pm\;6.4\times10^{-65}$ \\
\addlinespace[3pt]
$0.300$ & Our analytic & $0.39497522084295373429747165489072725825168951498178$ \\
 & Numerical difference & $2.72\times10^{-9}\;\pm\;6.8\times10^{-7}$ \\
 & Reference difference & $8.37\times10^{-90}\;\pm\;1.1\times10^{-69}$ \\
\addlinespace[3pt]
$4.50$ & Our analytic & $0.77515165401174580343102580969237768056470578000273$ \\
 & Numerical difference & $5.71\times10^{-8}\;\pm\;1.2\times10^{-6}$ \\
 & Reference difference & $2.40\times10^{-91}\;\pm\;3.7\times10^{-72}$ \\
\addlinespace[3pt]
$49.5$ & Our analytic & $0.35308838117224168652124506257313479270540259204215$ \\
 & Numerical difference & $-3.86\times10^{-8}\;\pm\;1.2\times10^{-6}$ \\
 & Reference difference & $-6.91\times10^{-89}\;\pm\;1.0\times10^{-71}$ \\
\addlinespace[3pt]
$1.00\times10^{3}$ & Our analytic & $0.070463896610137198990117049003399503906848383360028$ \\
 & Numerical difference & $3.33\times10^{-6}\;\pm\;1.5\times10^{-5}$ \\
 & Reference difference & $-4.10\times10^{-87}\;\pm\;6.7\times10^{-71}$ \\
\addlinespace[3pt]
$1.00\times10^{5}$ & Our analytic & $0.0028704528832006358222188994983077446681129388945786$ \\
 & Numerical difference & $-1.34\times10^{-6}\;\pm\;8.1\times10^{-6}$ \\
 & Reference difference & $-9.38\times10^{-87}\;\pm\;5.5\times10^{-70}$ \\
\bottomrule
\end{tabular}
\caption{Three independent evaluations of the NLO Compton correction
$\sigma_1=(\alpha_{\mathrm T}^3/m^2)C_1$, from low to high incident
energy in the electron rest frame. Our analytic coefficient $C_1$ is
evaluated from the collected formula and printed to 50 significant digits. Each difference subtracts this
coefficient from the indicated calculation. The $\pm$ entries estimate
numerical error, not experimental uncertainty. The reference is
Equation (6) of Lee, Schwartz and Zhang~\cite{qed-processes-lsz}.}
\label{tab:comp-independent-total-50}
\end{table}

